% ch07_u04 —— 原书 p.266--p.275（7.6.11 -- 7.8.6 前半）
% 本单元自 p.266 首句起（承 p.265 末句的谓词），至 p.275 末行止
% （该句跨页未完，尾部接 p.276）。

\SourcePage{266}

$\natural:\mathfrak{k}\to\operatorname{End}_{\mathcal{D}}M$ 是 $M$ 上的
$\mathfrak{k}$-作用。注意 $\natural$ 平凡，当且仅当 $M$ 是 $K$-等变
$\mathcal{D}$-模，即 $M\in\mathcal{M}(\mathcal{B})$。

\bdp{7.6.11} $Y$ 上的 $K$-等变 $\mathcal{D}$-复形，指由弱 $K$-等变
$\mathcal{D}$-模构成的复形 $N$，并带态射
$\mathfrak{k}\otimes N^{\cdot}\to N^{\cdot-1}$，$\xi\otimes n\mapsto i_\xi n$，
使得对任意 $\xi\in\mathfrak{k}$ 有 $i_\xi^2=0$，
$di_\xi+i_\xi d=\xi^{\natural}$。按记号滥用，把这种复形的 DG 范畴记作
$C(\mathcal{B},\mathcal{D})$。注意任意 $K$-等变 $\mathcal{D}$-模按显然方式
就是 $K$-等变 $\mathcal{D}$-复形；任意 $K$-等变 $\mathcal{D}$-复形的上同调层
都是 $K$-等变 $\mathcal{D}$-模。于是有上同调函子
$H:C(\mathcal{B},\mathcal{D})\to\mathcal{M}(\mathcal{B})$。把
$C(\mathcal{B},\mathcal{D})$ 的同伦范畴对 $H$-拟同构作局部化，得到三角范畴
$D(\mathcal{B},\mathcal{D})$。容易看出它是 $t$-范畴，其核心为
$\mathcal{M}(\mathcal{B})$。
% [Tutor Note: 原书 p.266 此处印作「such that for any xi in k our has」，
% 其中「our has」应为「one has」之误；底本如此，译文按「有」处理。]

对任意 $F\in C(\mathcal{B},\Omega)$，$\mathcal{D}$-复形 $\mathcal{D}F$
带上算子 $i_\xi^{\mathcal{D}F}=i_\xi^F\otimes\operatorname{id}_{\mathcal{D}Y}$，
是 $K$-等变的。对任意 $N\in C(\mathcal{B},\mathcal{D})$，$\Omega$-复形
$\Omega N$ 带算子 $i_\xi^{\Omega N}$；后者在
$N^i\otimes\Lambda^{-j}\Theta_Y$ 上按
$n\otimes\tau\mapsto i_\xi n\otimes\tau+(-1)^in\otimes\sigma(\xi)\wedge\tau$
作用，故 $\Omega N$ 是 $K_\Omega$-等变的 $\Omega$-复形。于是有伴随函子
$\mathcal{D}$、$\Omega$：
\begin{equation}\tag{322}
C(\mathcal{B},\Omega)\rightleftarrows C(\mathcal{B},\mathcal{D})
\end{equation}
以及三角范畴之间互逆的等价
\begin{equation}\tag{323}
D(\mathcal{B},\Omega)\rightleftarrows D(\mathcal{B},\mathcal{D}).
\end{equation}
后一个等价把 $D(\mathcal{B},\mathcal{D})$ 上上述 $t$-结构与 7.3.2 中定义的
$D(\mathcal{B},\Omega)$ 上的 $t$-结构等同起来。当所论的叠是光滑簇对群作用的商时，
这给出 7.3.4 的另一个证明。

\bdsec{7.7. Harish-Chandra 模及其导出范畴}

\SourcePage{267}

\bdp{7.7.1} 设 $G$ 是仿射代数群，$K\subset G$ 是代数子群，于是有
Harish-Chandra 对 $(\mathfrak{g},K)$。考虑范畴
$\mathcal{M}(K\backslash G/^{/}G)=\mathcal{M}((K\backslash G)/^{/}G)$，
其对象是 $G$ 上的 $\mathcal{D}$-模，带彼此交换的 $K$-作用与弱 $G$-作用
（其中 $K$ 与 $G$ 分别经左、右平移作用在 $G$ 上）。对
$M\in\mathcal{M}(K\backslash G/^{/}G)$，置
$\gamma(M)=\gamma^r(M):=\Gamma(G,M_G)^G$；这里把 $M_G$ 看作 $G$ 上的右
$\mathcal{D}$-模。这是 $(\mathfrak{g},K)$-模：$\mathfrak{g}$ 经右 $G$-平移
不变的向量场作用在 $\gamma(M)$ 上（按 $M$ 上的 $\mathcal{D}$-模结构），
$K$ 经左 $K$-平移作用。

\bdp{7.7.2} \emph{引理.} 函子
$\gamma:\mathcal{M}(K\backslash G/^{/}G)\to\mathcal{M}(\mathfrak{g},K)$
是范畴等价。

\emph{证明.} 留给读者（或见 [Kas]）。\hfill$\Box$

\bdp{7.7.3} \emph{注记.} (i) 置 $\gamma^l(M):=\Gamma(G,M^l_G)^G$，其中
$M^l_G$ 是 $M$ 的左 $\mathcal{D}$-模实现。由与上面相同的理由，这是
$(\mathfrak{g},K)$-模；有显然的等同
$\gamma^l(M)=\gamma^r(M)\otimes\det\mathfrak{g}$。

(ii) 有向量空间的典范同构 $\gamma^l(M)\cong M^l_{G,1}
=M^l_{K\backslash G,1}$，它把 $G$-不变截面送到该截面在 $1\in G$ 处的值。
$M^l_{K\backslash G,1}$ 上的 $(\mathfrak{g},K)$-模结构可描述如下。$K$-作用
来自右 $K$-平移在 $K\backslash G$ 上的（弱）作用（注意 $K$ 是
$1\in K\backslash G$ 的稳定化子），$\mathfrak{g}$-作用来自与弱 $G$-作用
相对应的 $\mathfrak{g}$ 的 $\natural$-作用（见 7.6.10）。

(iii) 设 $P$ 是 $K$-模，$\mathcal{P}$ 是 $K\backslash G$ 上相应的 $G$-等变
向量丛，其纤维为 $\mathcal{P}_1=P$。有
$\mathcal{D}\mathcal{P}=\mathcal{P}\otimes\mathcal{D}_{K\backslash G}
\in\mathcal{M}((K\backslash G)/^{/}G)$，且
$\gamma(\mathcal{D}\mathcal{P})=U(\mathfrak{g})\otimes_{U(\mathfrak{k})}
(P\otimes\det\mathfrak{k}^*)$。

\bdp{7.7.4} 上述引理给出 $(\mathfrak{g},K)$-模的导出范畴
$D(\mathfrak{g},K)$ 上典范的 $\mathcal{H}$-作用，这正是 7.1.1(c) 所许诺的。
事实上，由 7.6.8(v)（以及 7.6.9）知 $D(K\backslash G/^{/}G)$ 是
$\mathcal{H}$-模；而 7.7.2 把 $D(\mathfrak{g},K)$ 与这个范畴等同起来。
下面在 7.8.2 中给出该作用的另一种描述。它与当前定义的等价性在 7.8.9、
7.8.10(i) 中确立。

\SourcePage{268}

本节其余部分（7.7.5--7.7.11）是 Harish-Chandra 情形下
$\mathcal{D}$-$\Omega$ 等价的一个插叙；作为附带收获，7.7.12 给出
Bernstein--Lunts 定理 [BL]1.3 的一个简单证明。读者可以跳过这部分，
直接看 7.8。

\bdp{7.7.5} 下面是 7.7.2 对 $\Omega$-复形的版本。

设 $\Omega_{\mathfrak{g}}$ 是 $\mathfrak{g}$ 的 Chevalley 上链 DG-代数，故
$\Omega_{\mathfrak{g}}^{\cdot}=\Lambda^{\cdot}\mathfrak{g}^*$。它带
$K_\Omega$ 的典范「伴随」作用（记号见 7.6.3）。具体说，$K$ 以余伴随方式
作用在 $\Omega_{\mathfrak{g}}^{\cdot}$ 上；$\xi\in\mathfrak{k}
=\mathfrak{k}_\Omega^{-1}$ 作为 $\Omega_{\mathfrak{g}}^{\cdot}$ 的导子
$i_\xi$ 作用，把 $\nu\in\mathfrak{g}^*=\Omega_{\mathfrak{g}}^1$ 送到
$\langle\nu,\xi\rangle$。

$\Omega_{(\mathfrak{g},K)}$-\emph{复形}指一个 DG
$(\Omega_{\mathfrak{g}},K_\Omega)$-模；也就是说，它是一个复形，带
$\Omega_{\mathfrak{g}}$-作用与 $K_\Omega$-作用，二者关于 $K_\Omega$ 在
$\Omega_{\mathfrak{g}}$ 上的作用相容。对 $\Omega_{(\mathfrak{g},K)}$-复形
$T$，把 $\nu\in\mathfrak{g}^*=\Omega_{\mathfrak{g}}^1$、
$\xi\in\mathfrak{k}=\mathfrak{k}_\Omega^{-1}$ 在 $T$ 上的作用分别记作
$a_\nu$、$i_\xi$。把 $\Omega_{(\mathfrak{g},K)}$-复形的 DG 范畴记作
$C\Omega_{(\mathfrak{g},K)}$，其同伦范畴记作
$K\Omega_{(\mathfrak{g},K)}$。

对 $F\in C(K\backslash G/^{/}G,\Omega)$，置
$\gamma(F):=\Gamma(G,F_G)^G$。这是 $\Omega_{(\mathfrak{g},K)}$-复形。
事实上，$\Omega_{\mathfrak{g}}^{\cdot}$ 经与「$G$ 上对右 $G$-平移不变的
微分形式 DG 代数」的通常等同作用在它上面；而 $K_\Omega$ 作用在
$\gamma(F)$ 上，因为它作用在 $F_G$ 上（见 7.6.4、7.6.5）。

\bdp{7.7.6} \emph{引理.} 函子
$\gamma:C(K\backslash G/^{/}G,\Omega)\to C\Omega_{(\mathfrak{g},K)}$
是 DG 范畴的等价。

\emph{证明.} 留给读者。\hfill$\Box$

\bdp{7.7.7} 我们把 $(\mathfrak{g},K)$-复形与
$\Omega_{(\mathfrak{g},K)}$-复形等同于 $K\backslash G$ 上的弱 $G$-等变复形。
下面写出 Harish-Chandra 情形下的标准函子 $\mathcal{D}$ 与 $\Omega$。
方便起见，引入 DG Harish-Chandra 对
$(\mathfrak{k}_\Omega\times\mathfrak{g},K)$（结构嵌入
$\operatorname{Lie}K\hookrightarrow\mathfrak{k}_\Omega\times\mathfrak{g}$
是对角映射）。

设 $DR_{\mathfrak{g}}$ 是 $\mathfrak{g}$ 的、系数在 $U\mathfrak{g}$ 中的
Chevalley 上链复形（把 $U\mathfrak{g}$ 看作左 $U\mathfrak{g}$-模），故
$DR_{\mathfrak{g}}^i=\Lambda^i\mathfrak{g}^*\otimes U\mathfrak{g}$。
于是 $DR_{\mathfrak{g}}$ 是 $\Omega_{\mathfrak{g}}$-复形，也是
$(\mathfrak{k}_\Omega\times\mathfrak{g},K)$-复形；这两个作用相容。

\SourcePage{269}

（这里 $(\mathfrak{k}_\Omega\times\mathfrak{g},K)$ 经投影
$(\mathfrak{k}_\Omega\times\mathfrak{g},K)\to(\mathfrak{k}_\Omega,K)$
作用在 $\Omega_{\mathfrak{g}}$ 上，见 7.7.5）。具体说，对
$\nu\in\Omega_{\mathfrak{g}}$，
$\varepsilon=(\varepsilon_l,\varepsilon_r)\in\mathfrak{k}\times\mathfrak{g}
=\mathfrak{k}_\Omega^0\times\mathfrak{g}$，
$\xi\in\mathfrak{k}=\mathfrak{k}_\Omega^{-1}$，$k\in K$ 与
$a=\alpha\otimes v\in DR_{\mathfrak{g}}$，有
\[
\nu a=\nu\alpha\otimes v,\qquad
\varepsilon a=\operatorname{Ad}_{\varepsilon_l}(\alpha)\otimes v
+\alpha\otimes(\varepsilon_l v-v\varepsilon_r),
\]
\[
\xi a=i_\xi(\alpha)\otimes v,\qquad
ka=\operatorname{Ad}_k(\alpha)\otimes\operatorname{Ad}_k(v).
\]

对 $(\mathfrak{g},K)$-模的复形（简称 $(\mathfrak{g},K)$-复形）$V$，置
$\Omega V:=\operatorname{Hom}_{\mathfrak{g}}(DR_{\mathfrak{g}},V)$；
按显然方式这是 $\Omega_{(\mathfrak{g},K)}$-复形。对
$\Omega_{(\mathfrak{g},K)}$-复形 $T$，置
$\mathcal{D}T=\mathcal{D}_{(\mathfrak{g},K)}T:=
T\otimes_{\Omega_{\mathfrak{g}},\mathfrak{k}_\Omega}DR_{\mathfrak{g}}
=(T\otimes_{\Omega_{\mathfrak{g}}}DR_{\mathfrak{g}})_{\mathfrak{k}_\Omega}$；
这是 $(\mathfrak{g},K)$-复形。于是有伴随 DG 函子
\begin{equation}\tag{324}
\mathcal{D}=\mathcal{D}_{(\mathfrak{g},K)}:
C\Omega_{(\mathfrak{g},K)}\longrightarrow C(\mathfrak{g},K),\qquad
\Omega:C(\mathfrak{g},K)\longrightarrow C\Omega_{(\mathfrak{g},K)}.
\end{equation}

\emph{注.} 对上述 $T$，令
$\overline{T}^{\cdot}\subset T^{\cdot}$ 为所有算子 $i_\xi$
（$\xi\in\mathfrak{k}$）的核。这是 $T^{\cdot}$ 的 $K$-子模与
$\Lambda^{\cdot}(\mathfrak{g}/\mathfrak{k})^*$-子模（这里
$\Lambda^{\cdot}(\mathfrak{g}/\mathfrak{k})^*\subset
\Lambda^{\cdot}\mathfrak{g}^*=\Omega_{\mathfrak{g}}^{\cdot}$），
且显然的态射
\begin{equation}\tag{325}
\Omega_{\mathfrak{g}}^{\cdot}
\otimes_{\Lambda^{\cdot}(\mathfrak{g}/\mathfrak{k})^*}
\overline{T}^{\cdot}\longrightarrow T^{\cdot},\qquad
\overline{T}^{\cdot}\otimes_{U\mathfrak{k}}U\mathfrak{g}
\longrightarrow\mathcal{D}T^{\cdot}
\end{equation}
是同构。

\bdp{7.7.8} 回到几何情形。有显然的等同
$\Gamma(G,DR_G)^G=DR_{\mathfrak{g}}$（记号见 7.2.2；$G$ 经右平移作用在
自己上）。对 $M\in C((K\backslash G)/^{/}G,\mathcal{D})$，有
$\Omega_{(\mathfrak{g},K)}$-复形的典范同构
$\gamma(\Omega M)\cong\Omega(\gamma M)$，它定义为复合
\[
\Gamma(G,\mathit{Hom}_{\mathcal{D}_G}(DR_G,M_G))^G
=\operatorname{Hom}_{\mathcal{D}_G}(DR_G,M_G)^G
=\operatorname{Hom}_{U\mathfrak{g}}(DR_{\mathfrak{g}},\gamma M).
\]
对 $F\in C(K\backslash G/^{/}G,\Omega)$，有类似的典范同构
$\gamma\mathcal{D}F\cong\mathcal{D}\gamma F$，其定义留给读者。

\bdp{7.7.9} 对 $\Omega_{(\mathfrak{g},K)}$-复形 $T$，置
$H_{\mathfrak{g}}^{\cdot}T=H^{\cdot}\mathcal{D}T
\in\mathcal{M}(\mathfrak{g},K)$。于是
$H_{\mathfrak{g}}^{\cdot}:K\Omega_{(\mathfrak{g},K)}\to
\mathcal{M}(\mathfrak{g},K)$ 是上同调函子。把
$K\Omega_{(\mathfrak{g},K)}$ 中诱导 $H_{\mathfrak{g}}^{\cdot}$ 之间同构的
态射称为 $\mathfrak{g}$-\emph{拟同构}。$\mathfrak{g}$-拟同构构成局部化族；
把 $K\Omega_{(\mathfrak{g},K)}$ 的相应局部化记作
$D\Omega_{(\mathfrak{g},K)}$。函子 $\mathcal{D}$、$\Omega$ 给出互逆的
导出范畴等价

\SourcePage{270}

\begin{equation}\tag{326}
D\Omega_{(\mathfrak{g},K)}\rightleftarrows D(\mathfrak{g},K),
\end{equation}
其中 $D(\mathfrak{g},K):=D\mathcal{M}(\mathfrak{g},K)$。等价 $\gamma$
给出导出范畴的等价
\begin{equation}\tag{327}
D(K\backslash G/^{/}G,\Omega)\cong D\Omega_{(\mathfrak{g},K)},\qquad
D((K\backslash G)/^{/}G,\mathcal{D})\cong D(\mathfrak{g},K).
\end{equation}

\bdp{7.7.10} \emph{注记.} (i) 任意 $\mathfrak{g}$-拟同构都是拟同构；
反之未必。

(ii) 任意 $\Omega_{(\mathfrak{g},K)}$-复形 $T$ 都可看作
$\Omega_{\mathfrak{g}}=\Omega_{(\mathfrak{g},1)}$-复形（忘掉
$K_\Omega$-作用），于是有相应的 $\mathfrak{g}$-模复形
$\mathcal{D}_{\mathfrak{g}}T:=T\otimes_{\Omega_{\mathfrak{g}}}
DR_{\mathfrak{g}}$。显然的投影 $\mathcal{D}_{\mathfrak{g}}T\to
\mathcal{D}_{(\mathfrak{g},K)}T$ 是拟同构。这意味着：
$\Omega_{(\mathfrak{g},K)}$-复形的态射是 $\mathfrak{g}$-拟同构，
当且仅当它是 $\Omega_{\mathfrak{g}}$-复形的 $\mathfrak{g}$-拟同构。

\bdp{7.7.11} 7.7.7、7.7.9 的格式有如下版本。回忆 $DR_{\mathfrak{g}}$ 是
$(\mathfrak{k}_\Omega\times\mathfrak{g},K)$-复形。于是上述
$\mathcal{D}_{\mathfrak{g}}T$ 是
$(\mathfrak{k}_\Omega\times\mathfrak{g},K)$-复形；而对
$(\mathfrak{k}_\Omega\times\mathfrak{g},K)$-复形 $V$，复形
$\Omega V:=\operatorname{Hom}_{\mathfrak{g}}(DR_{\mathfrak{g}},V)$ 是
$\Omega_{(\mathfrak{g},K)}$-复形。函子
\begin{equation}\tag{328}
\mathcal{D}_{\mathfrak{g}}:C\Omega_{(\mathfrak{g},K)}\longrightarrow
C(\mathfrak{k}_\Omega\times\mathfrak{g},K),\qquad
\Omega:C(\mathfrak{k}_\Omega\times\mathfrak{g},K)\longrightarrow
C\Omega_{(\mathfrak{g},K)}
\end{equation}
是伴随的，同伦范畴之间的相应函子亦然。过渡到导出范畴，它们成为
（用 7.7.10(ii)）互逆的等价
\begin{equation}\tag{329}
D\Omega_{(\mathfrak{g},K)}\rightleftarrows
D(\mathfrak{k}_\Omega\times\mathfrak{g},K).
\end{equation}
投影 $(\mathfrak{k}_\Omega\times\mathfrak{g},K)\to(\mathfrak{g},K)$ 给出
全忠实嵌入 $C(\mathfrak{g},K)\to
C(\mathfrak{k}_\Omega\times\mathfrak{g},K)$，从而有正合函子
\begin{equation}\tag{330}
D(\mathfrak{g},K)\longrightarrow
D(\mathfrak{k}_\Omega\times\mathfrak{g},K).
\end{equation}
下述定理属于 Bernstein 与 Lunts [BL]1.3\footnote{[BL] 的作者只考虑有界导出范畴。}：

\SourcePage{271}

\bdp{7.7.12} \emph{定理.} 函子（330）是范畴等价。

\emph{证明.} (328) 中的函子 $\Omega$ 限制到 $C(\mathfrak{g},K)$ 上，
与 (324) 中的 $\Omega$ 一致。于是 7.7.12 由（326）与（329）推出。逆函子
$D(\mathfrak{k}_\Omega\times\mathfrak{g},K)\to D(\mathfrak{g},K)$ 把 $V$
送到 $\mathcal{D}_{(\mathfrak{g},K)}\Omega V$。\hfill$\Box$

\bdsec{7.8. Hecke 作用与局部化函子}

\bdp{7.8.1} 我们将描述 Harish-Chandra 模的导出范畴上典范的 Hecke 作用。
我们考虑扭情形，即 $\mathfrak{g}$ 的中心扩张的表示。以下是特征标的清单。

设 $G'$ 是 $G$ 被 $\mathbb{G}_m$ 的中心扩张，并带分裂 $K\to G'$。于是
$K$ 的原像 $K'\subset G'$ 与 $K\times\mathbb{G}_m$ 等同。置
$\mathfrak{g}':=\operatorname{Lie}G'$，
$\mathfrak{k}':=\operatorname{Lie}K'
=\mathfrak{k}\times\mathbb{C}$。有 Harish-Chandra 对
$(\mathfrak{g}',K')$，以及配套的 DG 对
$(\mathfrak{k}_\Omega\times\mathfrak{g}',K')$（这里结构嵌入
$\mathfrak{k}'\hookrightarrow
\mathfrak{k}_\Omega\times\mathfrak{g}'$ 的第一个分量是投影
$\mathfrak{k}'\to\mathfrak{k}$）。

设 $\mathcal{M}(\mathfrak{g},K)'$ 是由 $(\mathfrak{g}',K')$-模构成的
范畴，其中 $\mathbb{G}_m\subset K'$ 按标准特征标作用；其对象称为
$(\mathfrak{g},K)'$-模，或简称 Harish-Chandra 模。这是 Abel 范畴。
类似地，设 $C(\mathfrak{k}_\Omega\times\mathfrak{g},K)'$ 是由那些
$(\mathfrak{k}_\Omega\times\mathfrak{g}',K')$-复形构成的范畴，其中
$\mathbb{G}_m$ 按标准特征标作用；其对象称为
$(\mathfrak{k}_\Omega\times\mathfrak{g},K)'$-复形，或简称
Harish-Chandra 复形。这是 DG 范畴，带显然的上同调函子，取值于
$\mathcal{M}(\mathfrak{g},K)'$。把相应的导出范畴记作
$D(\mathfrak{g},K)'$；这是以 $\mathcal{M}(\mathfrak{g},K)'$ 为核心的
$t$-范畴。

\emph{注.} 由 Bernstein--Lunts 定理的扭版本，
$D(\mathfrak{g},K)'$ 等价于 $\mathcal{M}(\mathfrak{g},K)'$ 的导出范畴
\footnote{扭 Bernstein--Lunts 由非扭版本（见 7.7.12）施于 Harish-Chandra
对 $(\mathfrak{g}',K')$ 得到。}。下文不用这一事实，因为 Hecke 作用是
自然用 $(\mathfrak{k}_\Omega\times\mathfrak{g},K)'$-复形定义的。

\SourcePage{272}

\bdp{7.8.2} 现在定义 $D(\mathfrak{g},K)'$ 上典范的 $\mathcal{H}$-作用。
先定义预 Hecke 幺半 DG 范畴 $\mathcal{H}^c:=C(K\backslash G/K,\Omega)$
在 $C(\mathfrak{k}_\Omega\times\mathfrak{g},K)'$ 上的作用；Hecke 作用
在过渡到导出范畴之后才出现。

记 $\mathcal{L}_G$ 为 $G$ 上对应于 $\mathbb{G}_m$-torsor $G'\to G$ 的
线丛。$G$ 在自己上的左、右平移作用典范地提升为 $\mathcal{L}_G$ 上的
$G'$-作用。于是 $\mathcal{L}_G$ 的截面等同于 $G'$ 上的函数 $\phi$，
使得对 $c\in\mathbb{G}_m$、$g'\in G'$ 有
$\phi(cg')=c^{-1}\phi(g')$。因此 $\mathbb{G}_m\subset G'$ 经右平移在
$\mathcal{L}_G$ 的截面上作用，作用的特征标与标准特征标互逆。

取 Harish-Chandra 复形
$V\in C(\mathfrak{k}_\Omega\times\mathfrak{g},K)'$。置
$\mathcal{V}_G:=\mathcal{L}_G\otimes V$。则 $\mathcal{V}_G$ 是 $G$ 上
左 $\mathcal{D}$-模的复形。事实上，$\mathfrak{g}'$ 在 $\mathcal{L}_G$
上的无穷小右平移作用与 $V$ 上的 $\mathfrak{g}'$-作用张量起来，给出
$\mathcal{V}_G$ 上的 $\mathfrak{g}$-作用。$\mathcal{V}_G$ 上的左
$\mathcal{D}$-模结构使得左不变向量场经上述 $\mathfrak{g}$-作用作用在
$\mathcal{V}_G$ 上。$\mathcal{D}$-复形 $\mathcal{V}_G$ 关于左 $G'$-平移
弱等变：这些平移按 $\mathcal{L}_G$ 上的相应作用与 $V$ 上的平凡作用之
张量作用。因此由 7.6.10，它带典范的 $\mathfrak{g}'$-作用 $\natural$。

\emph{注.} 对 $\theta\in\mathfrak{g}'$，考虑函数
$\theta^{\natural}:G\to\mathfrak{g}'$，
$\theta^{\natural}(g):=\operatorname{Ad}_g(\theta)$。则对 $v\in V$、
$l\in\mathcal{L}_G$ 有
$\theta^{\natural}(l\otimes v)=l\otimes\theta^{\natural}(v)$。

取 $F\in\mathcal{H}^c$。则 $F_G\otimes\mathcal{V}_G$ 是 $G$ 上的
$\Omega$-复形（见 7.2.3(ii)）。它关于右 $K$-平移是 $K_\Omega$-等变的。
具体说，$K$ 按 $F$、$\mathcal{L}_G$ 上的相应作用与 $V$ 上的结构作用之
张量作用；算子 $i_\xi$ 按 $F$ 上右平移作用的相应算子与 $V$ 的结构算子
之和作用。把 $G/K$ 上相应的 $\Omega$-复形记作
$(F\otimes\mathcal{V})_{G/K}$。$F_G\otimes\mathcal{V}_G$ 上来自
$\mathcal{V}_G$ 上作用 $\natural$ 的 $\mathfrak{g}'$-作用与这个
$K_\Omega$-作用交换，故它定义了 $(F\otimes\mathcal{V})_{G/K}$ 上的
$\mathfrak{g}'$-作用。我们也把它记作 $\natural$。

\emph{注.} 若 $V$ 是 $(\mathfrak{g},K)'$-模的复形，则 $\mathcal{V}_G$ 是
左 $\mathcal{D}_G$-模的复形，关于右 $K$-平移强等变。

\SourcePage{273}

设 $\mathcal{V}_{G/K}$ 是 $G/K$ 上相应的左 $\mathcal{D}$-模复形。有
$(F\otimes\mathcal{V})_{G/K}=F_{G/K}\otimes\mathcal{V}_{G/K}$。

置 $F\overset{\sim}{\circledast}V:=
\Gamma(G,F_G\otimes\mathcal{V}_G)$，以及
\begin{equation}\tag{331}
F\overset{c}{\circledast}V=\Gamma(G/K,(F\otimes\mathcal{V})_{G/K})
=(F\overset{\sim}{\circledast}V)^{K_\Omega}.
\end{equation}

这是 $(\mathfrak{k}_\Omega\times\mathfrak{g},K)'$-复形。事实上，
$\mathfrak{g}'$ 按 $\natural$ 作用，$K$ 按 $F$ 与 $\mathcal{V}$ 的左平移
作用之张量作用，而算子 $i_\xi$ 是 $F$ 的相应算子。Harish-Chandra
相容性留给读者验证。

现在 $\overset{c}{\circledast}$ 定义了
$C(\mathfrak{k}_\Omega\times\mathfrak{g},K)'$ 上的
$\mathcal{H}^c$-模结构。事实上，结合约束
$(F_1\overset{c}{\circledast}F_2)\overset{c}{\circledast}V=
F_1\overset{c}{\circledast}(F_2\overset{c}{\circledast}V)$ 由显然的等同
\[
\Gamma(G,(F_1\overset{c}{\circledast}F_2)^{\cdot}\otimes\mathcal{L}_G)
=[\Gamma(G,F_1^{\cdot}\otimes\mathcal{L}_G)
\otimes\Gamma(G,F_2^{\cdot}\otimes\mathcal{L}_G)]^{K_\Omega}
\]
推出，其中 $K_\Omega$ 按右、左平移作用之张量作用（见 7.6.5）。我们把
Hecke 作用 $\circledast:\mathcal{H}\times D(\mathfrak{g},K)'\to
D(\mathfrak{g},K)'$ 定义为 $\overset{c}{\circledast}$ 的右导出函子。
若 $F$ 松，则 $F\circledast V=F\overset{c}{\circledast}V$，故
$\circledast$ 的结合约束由 $\overset{c}{\circledast}$ 的结合约束推出。

\emph{注.} 由上一则注，对
$M\in\mathcal{M}(K\backslash G/K)\subset\mathcal{H}$、
$V\in\mathcal{M}(\mathfrak{g},K)'$ 有
\begin{equation}\tag{332}
H^{\cdot}M\circledast V
=H_{DR}^{\cdot}(G/K,M\otimes\mathcal{V}_{G/K}).
\end{equation}

\bdp{7.8.3} \emph{注.} 设我们的扭是平凡的，故
$G'=G\times\mathbb{G}_m$。有显然的等价
$\mathcal{M}(\mathfrak{g},K)'=\mathcal{M}(\mathfrak{g},K)$ 与
$D(\mathfrak{g},K)=D(\mathfrak{g},K)'$（见 7.7.11）。于是我们定义了
$D(\mathfrak{g},K)$ 上的 Hecke 作用。7.8.9 将表明，这个作用确实与
7.7.4 中的作用一致。

回到一般情形。设 $U'$ 是 $\mathfrak{g}$ 的扭包络代数；用 $\mathfrak{z}$
记其由 $\operatorname{Ad}G$-不变元素构成的子代数。交换代数 $\mathfrak{z}$
按显然方式作用在任意 Harish-Chandra 复形上，故
$C(\mathfrak{k}_\Omega\times\mathfrak{g},K)'$ 从而 $D(\mathfrak{g},K)$
是 $\mathfrak{z}$-范畴。
% [Tutor Note: 原书 p.273 末句此处印作「hence D(g,K)」（g 为 fraktur 体），
% 无撇号；按上下文应为带撇的 D(g,K)。底本如此，保留原样。]

\SourcePage{274}

\bdp{7.8.4} \emph{引理.} $C(\mathfrak{k}_\Omega\times\mathfrak{g},K)'$、
$D(\mathfrak{g},K)'$ 上的 Hecke 作用是 $\mathfrak{z}$-线性的。

\emph{证明.} 用 7.8.2 中第一则注。\hfill$\Box$

\bdp{7.8.5} \emph{例.}（5 中将用到）令 $Vac':=U'/U'\mathfrak{k}$ 为扭
真空模。我们来显式计算 $F\circledast Vac'$。记号沿用 7.8.2。于是按 7.8.2
中第二则注，有 $G/K$ 上的左 $\mathcal{D}$-模 $\mathcal{V}_{G/K}$，它关于
左 $G$-平移弱等变，且 $\mathcal{V}_G=\mathcal{L}_G\otimes Vac'$。嵌入
$\mathbb{C}\subset Vac'$ 给出嵌入
$\mathcal{L}_{G/K}\subset\mathcal{V}_{G/K}$。容易看出，相应的左
$\mathcal{D}_{G/K}$-模态射 $\mathcal{D}_{G/K}
\otimes_{\mathcal{O}_{G/K}}\mathcal{L}_{G/K}\to\mathcal{V}_{G/K}$
是弱 $G$-等变 $\mathcal{D}$-模的同构。

\emph{注.} $\mathcal{D}_{G/K}\otimes\mathcal{L}_{G/K}$ 上与 $\natural$
对应的 $\mathfrak{g}'$-作用由下列公式给出：$\alpha'(\psi\otimes l)
=\psi\otimes\alpha'(l)-\psi\cdot\alpha\otimes l$，其中
$\alpha'\in\mathfrak{g}'$，$\alpha$ 是 $G/K$ 上相应的左平移向量场，
$\alpha'(l)$ 是 $l\in\mathcal{L}_{G/K}$ 的无穷小左平移。

于是对 $F\in\mathcal{H}^c$，有
$(F\otimes\mathcal{V})_{G/K}=F_{G/K}\otimes\mathcal{D}_{G/K}
\otimes_{\mathcal{O}_{G/K}}\mathcal{L}_{G/K}
=\mathcal{D}(F_{G/K})\otimes_{\mathcal{O}_{G/K}}\mathcal{L}_{G/K}$。
因此
\begin{equation}\tag{333}
F\overset{c}{\circledast}Vac'
=\Gamma(G/K,\mathcal{D}(F_{G/K})\otimes\mathcal{L}_{G/K}).
\end{equation}

这里 $\Gamma(G/K,\mathcal{D}(F_{G/K})\otimes\mathcal{L}_{G/K})$ 上的
$(\mathfrak{k}_\Omega\times\mathfrak{g},K)'$-作用定义如下。$\mathfrak{g}'$-
作用来自上面注中所述 $\mathcal{D}(F_{G/K})\otimes\mathcal{L}_{G/K}$ 上的
$\mathfrak{g}'$-作用；$K$-作用是左平移作用；算子 $i_\xi$ 来自
$F_{G/K}$ 上相应的算子。

过渡到导出函子（相当于在上述公式中考虑松的 $F$），得到
\begin{equation}\tag{334}
F\circledast Vac'=R\Gamma(G/K,
\mathcal{D}(F_{G/K})\otimes\mathcal{L}_{G/K}).
\end{equation}
特别地，对 $M\in\mathcal{M}(K\backslash G/K)$ 有
\begin{equation}\tag{335}
M\circledast Vac'=R\Gamma(G/K,
M_{G/K}\otimes\mathcal{L}_{G/K}).
\end{equation}

\SourcePage{275}

这里右端的 $\mathfrak{g}'$-作用来自
$M_{G/K}\otimes\mathcal{L}_{G/K}$ 上的 $\mathfrak{g}'$-作用，由公式
$\alpha'(m\otimes l)=m\otimes\alpha'(l)-m\alpha\otimes l$ 给出。

\bdp{7.8.6} 我们来解释 7.1.1 中「Hecke 模式」的 (d) 部分。先定义局部化
函子 $\Delta$。记号沿用 7.8.1。设 $Y$ 是 $G$ 作用其上的光滑簇，
$\mathcal{L}=\mathcal{L}_Y$ 是 $Y$ 上的线丛。设 $\mathcal{L}$ 带
$G'$-作用，该作用提升 $Y$ 上的 $G$-作用，并使
$\mathbb{G}_m\subset G'$ 按与标准特征标相反的特征标作用。线丛
$\omega_Y\otimes\mathcal{L}$ 带类似的作用。

我们如下定义一个 DG 函子
\begin{equation}\tag{336}
\Delta_\Omega=\Delta_{\Omega\mathcal{L}}:
C(\mathfrak{k}_\Omega\times\mathfrak{g},K)'\longrightarrow
C(K\backslash Y,\Omega)
\end{equation}
注意 $(\mathfrak{g}',K')$ 从而
$(\mathfrak{k}_\Omega\times\mathfrak{g}',K')$ 作用在
$\omega_Y\otimes\mathcal{L}$ 上（因为 $G'$ 如此）。对 Harish-Chandra
复形 $V$，考虑 $\mathcal{O}$-模的复形
$\omega_Y\otimes\mathcal{L}\otimes V$。
$\omega_Y\otimes\mathcal{L}$ 与 $V$ 上
$(\mathfrak{k}_\Omega\times\mathfrak{g}',K')$-作用的张量积给出
$\omega_Y\otimes\mathcal{L}\otimes V$ 上的
$(\mathfrak{k}_\Omega\times\mathfrak{g},K)$-作用。置
\[
\Delta_\Omega(V):=\operatorname{Hom}_{\mathfrak{g}}
(DR_{\mathfrak{g}},\omega_Y\otimes\mathcal{L}\otimes V)[-\dim K]
\]
（记号见 7.7.7）。换言之，$\Delta_\Omega(V)$ 是 $\mathfrak{g}$ 的、系数在
$\omega_Y\otimes\mathcal{L}\otimes V$ 中的平移后的 Chevalley 链复形。
这是 $Y$ 上的 $\Omega$-复形。由于 $DR_{\mathfrak{g}}$ 与
$\omega_Y\otimes\mathcal{L}\otimes V$ 都是
$(\mathfrak{k}_\Omega\times\mathfrak{g},K)$-复形，故
$\Delta_\Omega(V)$ 是 $K_\Omega$-等变的，即
$\Delta_\Omega(V)\in C(K\backslash Y,\Omega)$。

注意 $\Delta_\Omega(V)$ 带典范的递增有限滤过，其相邻商等于
$\Lambda^i\mathfrak{g}\otimes\omega_Y\otimes\mathcal{L}\otimes V
[i-\dim K]$。因此 $\Delta_\Omega$ 把拟同构送到 $\mathcal{D}$-拟同构。
于是它给出三角函子
\begin{equation}\tag{337}
L\Delta=L\Delta_{\mathcal{L}}:
D(\mathfrak{g},K)'\longrightarrow D(K\backslash Y)
\end{equation}
上述注还表明，$L\Delta$ 是右 $t$-正合函子。核心之间相应的右正合函子是
$\Delta_{\mathcal{L}}:\mathcal{M}(\mathfrak{g},K)'\to$
% 尾部接 p.276（原书此句跨页，p.275 末行止于箭头）。
